\section{Joint learning and finite certificates for a finite alphabet}
\label{sec:multilearning82}

We next construct one learner common to all unknown elements of
$\mathfrak P_{d,k}$.  Its training experiments do not use the
pair-dependent eigenvectors in the geometric lower bound.
The outcome count $k$ is public, as are $d,N,\delta,\eta$.
The upper procedure uses independent one-call probe--reference pairs;
coherent collective processing is confined to completed outputs.
The converse permits arbitrary coherent adaptive training, with an
every-record budget and bounded public stopping.

\begin{theorem}[Joint finite-outcome learning and coding]
\label{thm:multioutcomelearning82}
Let $d,N\ge1$, $k\ge2$, $0<\delta\le2^{-24}$, and
$0<\eta\le1/8$.  Write $M^*_{d,k}(N,\delta,\eta)$ for the
minimum every-record call budget with worst-case failure at most
$\eta$ under $D_N$-error $\delta$ on $\mathfrak P_{d,k}$.
There are absolute constants $c,C>0$ such that
\begin{equation}\label{eq:multilearning82}
 c\frac{N}{\delta^2}\left(d^2+d\log\frac1\eta\right)
 \le M^*_{d,k}(N,\delta,\eta)
 \le C\frac{k^3N}{\delta^2}
                  \left(d^2+d\log\frac{k}{\eta}\right).
\end{equation}
For each fixed $k$, this is a matching order jointly in
$d,N,\delta,\eta$, including $d=1$:
\[
 M^*_{d,k}(N,\delta,\eta)
 =\Theta_k\!\left(N\delta^{-2}
                   [d^2+d\log(1/\eta)]\right).
\]
The displayed dependence on growing $k$ is not claimed optimal.

For rational accuracy and confidence, the upper procedure has a
finite deterministic specification, returns an index with a fixed
public decoder into $\mathfrak P_{d,k}$, and simultaneously uses
\begin{equation}\label{eq:multilearnedbits82}
 B=\frac{(k-1)d^2}{2}\log_2N
       +(k-1)d^2\log_2(1/\delta)+O((k-1)d^2\log(k+1))
\end{equation}
bits.  On its good event the decoded error is at most $5\delta/8$.
The same call order is available with finite trusted controls.
Classical search, readout storage, known-state preparation, and
known-control realization are separate resources.
\end{theorem}
\begin{proof}
We first give the constructive upper bound for rational parameters.
Set $n=\lceil\sqrt N\rceil$ and
\begin{equation}\label{eq:multitrainingaccuracy82}
 \epsilon=\frac{\delta}{16kn}.
\end{equation}
For each $j$ consider the binary effect
$F_j=I_d/4+E_j/2$.  It belongs to $[I_d/4,3I_d/4]$.
One call to $\mathcal M_{\boldsymbol E}$ simulates a call to
$\mathcal M_{F_j}$: conditional on the original label $x$, return
bit one with probability $3/4$ when $x=j$ and with probability
$1/4$ otherwise.  This is an equality of channels, including their
action on a retained reference, because the bit-one effect is
\[
 \tfrac34E_j+\tfrac14\sum_{x\ne j}E_x=\tfrac14I_d+\tfrac12E_j.
\]
Only public classical randomness is added.  Each simulated call
uses exactly one call to the unknown device.

Apply Theorem~\ref{thm:rationallearner81} to each $F_j$ at future
horizon one, accuracy $\epsilon$, and failure $\eta/k$, using
fresh records for different $j$.  Its output $\widehat F_j$ is
Gaussian rational and legal.  The binary identity
$d_1(F_j,\widehat F_j)=2\|F_j-\widehat F_j\|_{\rm op}$ implies
that, with probability at least $1-\eta$, all the Hermitian
matrices $B_j=2\widehat F_j-I_d/2$ obey
\begin{equation}\label{eq:multiassembly82}
 \max_j\|B_j-E_j\|_{\rm op}\le\epsilon.
\end{equation}
Only a union bound is needed; no independence assertion concerning
selected quantum readout outcomes is used.

The matrices $B_j$ need not sum to $I_d$.  Reconstruct the finite
legal grid $\mathcal G(\epsilon/2)$ and choose its first tuple
$\boldsymbol G$ satisfying
\[
 \|G_j-B_j\|_{\rm op}\le2\epsilon\quad\text{for every }j.
\]
Each comparison is an exact rational positive-semidefinite test.
On \eqref{eq:multiassembly82}, such a tuple exists by
Lemma~\ref{lem:multigrid82}; it has
$r(\boldsymbol G,\boldsymbol E)\le3\epsilon$.
On every other record, if the finite search finds no tuple, return
$\boldsymbol U$.  Thus every record has a legal output and the
same predetermined training budget.  The geometric upper bound gives
\[
 D_N(\boldsymbol E,\boldsymbol G)
       \le6k\sqrt N\epsilon\le3\delta/8.
\]
Encode $\boldsymbol G$ with the dictionary of
Theorem~\ref{thm:multioutcomecode82} at parameter $\delta/2$.
Its additional error is at most $\delta/4$, so the decoded error
is at most $5\delta/8$.  It uses no further device calls and has
length \eqref{eq:multilearnedbits82}.

There are $k$ binary learners, each using at most
$C_0\epsilon^{-2}[d^2+d\log(k/\eta)]$ calls.  Since $n^2\le4N$,
this proves the upper bound.  The inherited rational synthesis
and finite-control theorem applies to the simulated binary channels
as to any channels of the same form.  The public $1/4,3/4$ relabeling
is exact classical control.  The error guarantee used in
\eqref{eq:multiassembly82} is the inherited actual finite-control
guarantee, not an idealized readout subsequently assumed exact.
The construction terminates in finite classical time, with no
assertion that this time is polynomial.  For nonrational error
parameters, choose rational $\delta'\in[\delta/2,\delta]$ and
$\eta'\in[\eta/2,\eta]$ to obtain the stated existence bound;
these choices change only absolute constants.

For the lower bound let $m=\lfloor k/2\rfloor$ and $q=2m/k\ge2/3$.
Embed any binary interior effect $F\in[I_d/4,3I_d/4]$ by the tuple
\begin{equation}\label{eq:multiembedding82}
 E_j(F)=\begin{cases}
             2F/k,&1\le j\le m,\\
             2(I_d-F)/k,&m<j\le2m,\\
             I_d/k,&j=k=2m+1.
             \end{cases}
\end{equation}
This belongs to $\mathfrak P_{d,k}$.  A query to this tuple is
simulated using one binary $F$ query: with probability $q$, retain
the binary outcome and select a uniform label from its corresponding
group of $m$ labels; otherwise erase the bit and return the extra
label.  In the odd case the erased reference state is the sum of
its two subnormalized binary states, as required by the effect
$I_d/k$.  Thus the simulation is valid for entangled inputs and
adaptive controls, not merely for outcome probabilities on pure
states.  For even $k$ there is no erasure.

Suppose a $k$-outcome learner returns any legal tuple
$\widehat{\boldsymbol E}$ within $D_N$-error $\delta$.
Coarse-graining the first $m$ labels gives a binary effect
$H=\sum_{j\le m}\widehat E_j$ satisfying
$d_N(qF,H)\le\delta$.  This follows by executing the same
coarse-graining inside any binary tester of the two induced channels.
The Bernoulli lower bound in
Theorem~\ref{thm:multioutcomemetric82}, applied to these two binary
effects, implies
\[
 \|H-qF\|_{\rm op}\le128\delta/\sqrt N,
\]
since $\delta<1/128$.
Spectrally clip $H/q$ to $[I_d/4,3I_d/4]$, obtaining $\widehat F$.
Weyl's inequality bounds the clipping displacement by
$\|H/q-F\|_{\rm op}$; the triangle inequality therefore gives
\[
 \|\widehat F-F\|_{\rm op}
      \le\frac{256}{q}\frac\delta{\sqrt N}
      \le\frac{384\delta}{\sqrt N}.
\]
No operator-Lipschitz claim about spectral clipping is used.
The last accuracy is at most $2^{-14}$ under the stated cap.
The proof of Corollary~\ref{cor:operatorconfidence80} gives the
binary operator-loss lower bound already on this interior, with
arbitrary coherent adaptive training.  Applying it to the simulation
proves the left inequality of \eqref{eq:multilearning82}, including
the contribution $d\log(1/\eta)$ and the scalar case.  The same
simulation preserves every-record budgets and public stopping.
\end{proof}

The length in \eqref{eq:multilearnedbits82} has the optimal leading
order also among learners with a fixed-length, fixed public decoder.
Indeed, if failure is less than one at every target, at least one
word must decode inside that target's loss ball.  The decoder's
image is therefore a covering of the entire target family.
The lower bound in Theorem~\ref{thm:multioutcomeentropy82} applies
regardless of the number of training calls.  This is separate from
the statistical query lower bound.  Shared-randomness expected-prefix
conventions are different and are not part of this assertion.

\subsection{Finite uniform-risk certificates}
The certification mechanism also survives a finite outcome alphabet.
For an ideal one-call maximally entangled acquisition, its normalized
classical--quantum record is
\begin{equation}\label{eq:multichoi82}
 \Omega(\boldsymbol E)=\frac1d\bigoplus_{j=1}^k E_j^{\mathsf T}.
\end{equation}
The classical labels remain ordered.  A collective readout after $M$
independent acquisitions can be described by one POVM
$(A_{c,x})_c$ on dimension $d^M$ for each classical string
$x\in\{1,\ldots,k\}^M$.  Here $A_{c,x}\succeq0$ and
$\sum_c A_{c,x}=I_{d^M}$.  Its output label $c$ is decoded by a
fixed finite family $\boldsymbol C_c$ of legal $k$-outcome tuples.

\begin{theorem}[Finite-alphabet certificate transfer]
\label{thm:multicertificate82}
Let $\mathcal G\subset\mathfrak P_{d,k}$ be a complete finite
$r_0$-net in the tuple norm $r$.  Suppose a fixed readout satisfies
\begin{equation}\label{eq:multicertificate82}
 \sum_{c:\,r(\boldsymbol G,\boldsymbol C_c)\le a}
 \sum_{x\in\{1,\ldots,k\}^M}
 \operatorname{tr}\left[A_{c,x}
        d^{-M}\bigotimes_{i=1}^M G_{x_i}^{\mathsf T}\right]
 \ge1-\alpha\qquad(\boldsymbol G\in\mathcal G).
\end{equation}
Then every real $\boldsymbol E\in\mathfrak P_{d,k}$ satisfies
\begin{equation}\label{eq:multitransfer82}
 \mathbb P_{\boldsymbol E}
   \{r(\boldsymbol E,\boldsymbol C_J)>a+r_0\}
 \le\alpha+\frac{Mk r_0}{2}.
\end{equation}
If all decoding centres are in $\mathfrak P_{d,k}$, the same failure
bound holds for $D_N$-error greater than $2k\sqrt N(a+r_0)$.
For Gaussian-rational data, rational $a,r_0,\alpha$, and the net
of Lemma~\ref{lem:multigrid82}, the hypotheses admit a finite exact
verification.  A proper prefix of the net or of the classical
string list is not a certificate.
\end{theorem}
\begin{proof}
The block-diagonal formula \eqref{eq:multichoi82} gives
\[
 \|\Omega(\boldsymbol E)-\Omega(\boldsymbol G)\|_1
    =d^{-1}\sum_j\|E_j-G_j\|_1
    \le k\,r(\boldsymbol E,\boldsymbol G).
\]
Tensor telescoping between normalized states gives
\begin{equation}\label{eq:multirecordcontinuity82}
 \|\Omega(\boldsymbol E)^{\otimes M}
           -\Omega(\boldsymbol G)^{\otimes M}\|_1
       \le Mk\,r(\boldsymbol E,\boldsymbol G).
\end{equation}
Apply the same readout to both states.  Its index laws differ in
total variation by at most $Mk r_0/2$ when
$r(\boldsymbol E,\boldsymbol G)\le r_0$.
Choose such a $\boldsymbol G$.  The fixed set of labels
$\{c:r(\boldsymbol G,\boldsymbol C_c)\le a\}$ has probability
at least $1-\alpha$ at $\boldsymbol G$, by
\eqref{eq:multicertificate82}; every label in this set has
$r(\boldsymbol E,\boldsymbol C_c)\le a+r_0$.
Transferring this one event proves \eqref{eq:multitransfer82}.
There is no assumption that the success-label sets vary continuously
with the unknown target.

The geometric upper comparison proves the assertion about $D_N$.
For exact verification, reconstruct the whole legal net and every
classical string.  All traces in \eqref{eq:multicertificate82} are
Gaussian-rational expressions with rational real values.  Check
complete POVM positivity and normalization at every string, use
$sI_d\pm(G_j-C_{c,j})\succeq0$ to determine the good labels
including equality, and check each rational inequality.  These are
finite computations.  Their soundness on the real parameter body
is supplied by \eqref{eq:multirecordcontinuity82}, not by a numerical
sampling argument.
\end{proof}

The certificate describes the ideal specified rational readout.
When implemented controls have a separately certified output-law
error $\zeta$, the right side of \eqref{eq:multitransfer82} increases
by at most $\zeta$.  Such an error must be budgeted, not silently
identified with zero.  The finite-control guarantee in
Theorem~\ref{thm:multioutcomelearning82} instead uses the already
budgeted binary procedures.  The exact certificate transfer and
the actual-control learning guarantee thus have explicitly distinct
hypotheses.
